%% kapa: exam=1 round=20 year=2009 period=1 form=A subject=economics subject_ko=경제원론 seq=39 number=79
%% source: Q-Net 감정평가사 기출문제 (cst00309) artlSeq=5024178 file=extracted/1008/5024178/2009년 제20회 감정평가사 1차시험 기출문제/1교시.hwp
%% conversion: kordoc 4.24.51 -> tools/gen1.py (auto); review=proofread-sonnet; flags=none
\begin{question}{79}
소비자 A는 1기와 2기에 걸쳐 소비를 한다. $c_{ 1 }$ 을 1기의 소비, $c_{ 2 }$ 를 2기의 소비라고 할 때, 소비자 A의 효용함수는 $U (c_{ 1 } , c_{ 2 } )=min \left \{ c_{ 1 } , c_{ 2 } \right \}$이다. 1기의 소득은 210만원이고, 2기의 소득은 0원이며, 각 기의 소비재 가격은 1원으로 동일하다. A는 1기에 10\%의 이자율로 저축을 하거나 대출을 받을 수 있다. 소비자 A의 행동 중 합리적인 것은?
\begin{choices}
  \item 100만원을 저축한다.
  \item 110만원을 저축한다.
  \item 100만원을 대출받는다.
  \item 110만원을 대출받는다.
  \item 저축하지도 대출받지도 않는다.
\end{choices}
\end{question}
\answer{1}
